% The print PDF's preamble (memoir, A5, lettrine drop caps, running heads, \Liber, \Caput, \Versus): the
% preamble of the former LaTeX/main.tex, moved here 2026-11-11 (decision 13). pipeline/pdf.py reads it, inlines
% title.tex, and adds its own commands; the body comes from the USFM master (latex.to_latex).
\documentclass[9pt,a5paper,twoside,openany]{memoir}
\usepackage[utf8]{inputenc}
\usepackage{layout}
\usepackage{lettrine}
\usepackage{float}
\usepackage{hyperref}
% multicol before footmisc: footmisc sets its run-in notes across the whole page only if multicol is loaded when it
% is; otherwise each page's notes took one column's width, except on a book's last page (2026-11-12).
\usepackage{multicol,microtype,etoolbox}
\usepackage[hang,flushmargin,para,perpage]{footmisc}
% footmisc reckons a note's height from \columnwidth, a column's inside multicols: by \textwidth, the width the notes
% are set to, so it reserves the lines they take, not twice as many.
\makeatletter
% That height is then rounded up to whole lines of the text, and the space above the notes is one line: TeX's page
% builder fills the two columns as one long one, but each column holds only whole lines, so a note reserving part
% of a line let it take one line more than the columns could; multicol carried it over to the next page and the
% note stayed behind (Lev 21:4's note at the foot of the page before the verse; found 2026-11-16).
\newdimen\CastellioLine
\AtBeginDocument{\CastellioLine=\baselineskip \skip\footins=\CastellioLine}
% Half a note line is added first: the notes, run together, end in a part line and begin with \footnotesep.
\def\Castellio@wholelines{\count@\ht\FN@tempboxa \advance\count@\dimexpr.5\footnotebaselineskip\relax
    \advance\count@\CastellioLine \advance\count@\m@ne
    \divide\count@\CastellioLine \ht\FN@tempboxa\count@\CastellioLine}
\patchcmd{\FN@footnotetext}{\footnotebaselineskip / \columnwidth\relax}
  {\footnotebaselineskip / \textwidth\relax \Castellio@wholelines}{}
  {\PackageWarning{castellio}{footmisc's note height not patched}}
% footmisc copied \FN@footnotetext as it loaded (into hyperref's \H@@footnotetext, as hyperref is loaded), so the
% patch above reached no note until the copy was made again (2026-11-16; the 2026-11-12 patch had never applied).
\@ifpackageloaded{hyperref}{\let\H@@footnotetext\FN@footnotetext}{\let\@footnotetext\FN@footnotetext}
% footmisc joins a page's notes into one paragraph in LaTeX's own \@makecol; memoir replaces \@makecol, so the notes
% stayed as one unbroken line each, and a long one ran off the page (Lev 13:48, 21:4; found 2026-11-16). memoir's
% footnote step now makes the paragraph first. (memoir's own \paragraphfootnotes uses an insertion multicol leaves
% out of its reckoning, so the columns could run into the notes.)
% An empty note box (below) is voided, not made a paragraph: that paragraph was a bare 5.6pt strut, overfilling the page.
\def\Castellio@notes{\ifdim\dimexpr\ht\footins+\dp\footins\relax>\z@
    \global\setbox\footins\vbox{\FN@makefootnoteparagraph}\else\global\setbox\footins\box\voidb@x\fi}
\pretocmd{\m@mopfootnote}{\Castellio@notes}{}
  {\PackageWarning{castellio}{memoir's footnotes not made a paragraph}}
\pretocmd{\m@mopfootnotebf}{\Castellio@notes}{}
  {\PackageWarning{castellio}{memoir's footnotes not made a paragraph}}
% When a book's columns end by filling a page that has notes, multicol puts an empty note (its placeholder) after
% them, which falls on the next page; \clearpage then made a page of it, blank but for its head (after Romans,
% whose last page has the note on 15:29; Nathan, 2026-11-20). \clearpage now voids an empty note box first.
\pretocmd{\@doclearpage}{\ifvoid\footins\else\ifdim\dimexpr\ht\footins+\dp\footins\relax=\z@
    \global\setbox\footins\box\voidb@x\fi\fi}{}{\PackageWarning{castellio}{\string\clearpage\space not patched}}
\makeatother
\usepackage{graphicx}   % \scalebox, for \TitulusMinor

\usepackage[LGR,T1]{fontenc}
% Greek in the footnotes (decision 53): latex.note_to_latex writes it as Unicode in \Graece{…}. Under pdfLaTeX it
% is set in the LGR encoding's TeX Greek fonts (greek-inputenc maps the Unicode); under LuaLaTeX or XeTeX
% pipeline/pdf.py redefines \Graece to GFS Didot, the Greek drawn to sit with Palatino.
\newcommand{\Graece}[1]{{\fontencoding{LGR}\selectfont #1}}
% Hebrew in the footnotes (Exod 40:20): set by LuaLaTeX or XeTeX (pipeline/pdf.py); pdfLaTeX has no Hebrew font, so
% there the word is left out.
\newcommand{\Hebraice}[1]{}
% \usepackage{imfellEnglish}
\usepackage{palatino}

\setlrmarginsandblock{0.5in}{0.6in}{*}
\setulmarginsandblock{0.8in}{0.6in}{*}

% The line spacing is set before the layout is fixed, so the text block holds a whole number of lines (\topskip +
% n \baselineskip): otherwise no column can be filled exactly, and multicols ends columns early.
\linespread{1.1}\selectfont
\checkandfixthelayout[lines]

\frenchspacing
% No paragraph ends on the stub of a hyphenated word ("Abrahā- / mō:", Matt 1:1).
\finalhyphendemerits=100000
% No stretch between paragraphs: a column filled out to full height must not open a gap where a paragraph meets
% a poem (Matt 1:1–2).
\setlength{\parskip}{0pt}
% No widows: a paragraph's last line never stands alone at the top of a column (Gen 25:33's "inquit" overleaf;
% Nathan, 2026-11-16). A column ends a line short instead, as at a chapter's opening.
\widowpenalty=10000
% \sloppy

\renewcommand*\footnoterule{}

\newcounter{libercount}
\newcounter{caputcount}[libercount]
\newcounter{versuscount}[caputcount]

\newcommand{\libernameshort}{}
\newcommand{\libernamelong}{}

% A recto page: a blank verso before it, if needed, carries no running head.
\makeatletter
\newcommand{\ClearToRecto}{\clearpage\ifodd\c@page\else\thispagestyle{empty}\null\clearpage\fi}
\makeatother

% Each book is set in two balanced columns (multicols), its title across both. The contents entry and the PDF
% bookmark are made after the page break, so they point to the book's first page; the first book opens recto.
% \FinisLibri (written by pipeline/pdf.py after each book) closes the columns, balancing the last page.
\newcommand{\Liber}[2][Missing Longform Title]{%
    \ifnum\value{libercount}=0 \ClearToRecto\else\clearpage\fi
    \renewcommand{\libernamelong}{#1}%
    \renewcommand{\libernameshort}{#2}%
    \stepcounter{libercount}%
    \phantomsection
    \addcontentsline{toc}{chapter}{#2}%
    \thispagestyle{empty}%
    \begin{multicols}{2}[\LiberTitulus]%
}
% A book title's \mt2 line, smaller, inside \LiberTitulus's \HUGE: scaled as a box, not by a size command. A \Large
% (or \fontsize) here made lettrine set every chapter's drop cap at 36.5pt, not 21.8pt, rising above its two lines
% (found 2026-11-12 by compiling the whole Bible with Tectonic; a single book was not affected).
% The whole title is scaled up 1.2 times (no size command larger than \HUGE, as above), in a box as wide as the text
% once scaled (10/12); \TitulusMinor is .55/1.2, so the \mt2 line keeps its size (2026-11-12: \HUGE small caps alone
% looked undersized).
\newcommand{\TitulusMinor}[1]{\scalebox{.46}{#1}}
% The break before a minor line under a full-size one (James, John's and Peter's letters: "Jācōbī Apostolī" over
% "Epistula"): the lines stand a \HUGE line apart, which under a minor line above leaves the gap between its foot and
% the capitals below, but over one adds the capitals' missing height (1 − .46 of it), so the gap was twice as wide
% (Nathan, 2026-11-20). That much is taken back. (A macro, so no [ ] stands inside \Liber's optional argument.)
\newcommand{\TitulusInfra}{\\[-.54\fontcharht\font`\E]}
\newcommand{\LiberTitulus}{%            % a macro: multicols' [title] argument may not hold a \par itself
    \begin{center}
        \scalebox{1.2}{\begin{minipage}{\dimexpr\textwidth*10/12\relax}
            \centering\HUGE \textsc{\libernamelong}\par
        \end{minipage}}\par
    \end{center}
    \vspace{1.5\baselineskip}}
\newcommand{\FinisLibri}{\par\end{multicols}}

% A chapter: one blank line before it (none at a book's start), then the drop numeral. Its first paragraph never
% leaves fewer than three lines at the foot of a column (\clubpenalties, for that paragraph only); a column may
% end up to two lines short instead. (Not \needspace, which measures the page and so breaks multicols.)
\newcommand{\Caput}{%
    \stepcounter{caputcount}%
    \par
    \ifnum\value{caputcount}>1
        \vskip 0pt plus 2\baselineskip \penalty0 \vskip 0pt plus -2\baselineskip
        \vspace{\baselineskip}\fi
    \clubpenalties 3 10000 10000 \clubpenalty
    \AddToHookNext{para/after}{\clubpenalties 0 }%
    \lettrine[findent=.4em,nindent=0pt]{\thecaputcount}{}%
}
% The numeral's gap: every line beside it starts .4em after it (findent; nindent 0, so the first line starts where
% the second does). The text is not lettrine's argument: its first \Versus (verse 1) prints no space, which as
% interword glue had stretched with the line, so the first line's start varied (Nathan, 2026-11-16).
% The numeral's top: level with the first line's capitals. lettrine aims at the height of its \LettrineSecondString,
% "x" by default (for an opening set in small caps), which left the numeral at the x-height (Nathan, 2026-11-16).
\renewcommand*{\LettrineSecondString}{E}

% A verse number; each one also sets a mark {book}{chapter}{verse} for the running heads, in a mark class of its own
% (LaTeX's \InsertMark). multicol keeps such classes to the lines actually set on the page; TeX's own \botmark also
% saw the lines multicol gathered for the page and carried over to the next, so a head could name a verse that
% begins overleaf ("Lēvīticus 22:14 – 23:15", 23:15 on the next page; found 2026-11-16).
\makeatletter
\NewMarkClass{versus}
% A verse that opens a paragraph starts the paragraph before its mark (\leavevmode): a mark set between paragraphs
% stood after the last line of the one before, so when the page broke there it went with that line, and both heads
% were a verse off (Gen 6:5 heading the page before its own, 6:6 its own; Nathan, 2026-11-20).
\newcommand{\Versus}{%
    \leavevmode
    \stepcounter{versuscount}%
    \protected@edef\@tempa{\noexpand\CastellioRef{\libernameshort}{\arabic{caputcount}}{\arabic{versuscount}}}%
    \InsertMark{versus}{\unexpanded\expandafter{\@tempa}}%   % \InsertMark expands its text: keep \CastellioRef
    \ifnum\value{versuscount}=1 \expandafter\ignorespaces\else
      {\textsuperscript{\theversuscount}}\nobreak\hskip\z@skip\fi%
}
% (The zero glue after the number lets TeX hyphenate the verse's first word: it hyphenates only a word that follows
% glue, and Heb 13:2's *Hospitālitātem*, unbroken, ran 10pt into the margin; 2026-11-16.)
% A verse that stands for a range of the KJV's (\v 12-83: Num 7, the princes' offering given once): numbered
% "12–83", the running head marking its first verse; the counter goes on from the last.
\newcommand{\VersusUsque}[1]{%
    \leavevmode
    \stepcounter{versuscount}%
    \protected@edef\@tempa{\noexpand\CastellioRef{\libernameshort}{\arabic{caputcount}}{\arabic{versuscount}}}%
    \InsertMark{versus}{\unexpanded\expandafter{\@tempa}}%
    \textsuperscript{\theversuscount--#1}\nobreak\hskip\z@skip
    \setcounter{versuscount}{#1}%
}
% The range of verses that begin on the page: "Genesis 1:31 – 2:25", "Genesis 3:1–14", "Genesis 3:15". A page
% with no verse mark yet in the document has an empty head. \CastellioRef is \relax (unexpandable) except while a
% head reads the marks: multicol expands them again when it carries them over.
\let\CastellioRef\relax
\def\Castellio@first#1#2#3{\def\Castellio@book{#1}\def\Castellio@fc{#2}\def\Castellio@fv{#3}}
\def\Castellio@last#1#2#3{\def\Castellio@lc{#2}\def\Castellio@lv{#3}}
\newcommand{\VersusRange}{{%
    \edef\@tempb{\LastMark{versus}}%
    \ifx\@tempb\@empty\else
        \let\CastellioRef\Castellio@first \FirstMark{versus}%
        \let\CastellioRef\Castellio@last \LastMark{versus}%
        \Castellio@book\ \Castellio@fc:\Castellio@fv
        \ifnum\Castellio@fc=\Castellio@lc
          \ifnum\Castellio@fv=\Castellio@lv\else\textendash\Castellio@lv\fi
        \else\ \textendash\ \Castellio@lc:\Castellio@lv\fi
    \fi}}
% The heads as most Bibles set them (Nathan, 2026-11-16): a verso names the first verse that begins on it, a recto
% the last, so an open spread reads from its left head to its right one ("Genesis 9:27" … "Genesis 11:4").
% \VersusRange is kept for any page style that wants the whole range.
\def\Castellio@lastbook#1#2#3{\def\Castellio@book{#1}\def\Castellio@lc{#2}\def\Castellio@lv{#3}}
\newcommand{\VersusPrimus}{{%
    \edef\@tempb{\FirstMark{versus}}%
    \ifx\@tempb\@empty\else
        \let\CastellioRef\Castellio@first \FirstMark{versus}%
        \Castellio@book\ \Castellio@fc:\Castellio@fv
    \fi}}
\newcommand{\VersusUltimus}{{%
    \edef\@tempb{\LastMark{versus}}%
    \ifx\@tempb\@empty\else
        \let\CastellioRef\Castellio@lastbook \LastMark{versus}%
        \Castellio@book\ \Castellio@lc:\Castellio@lv
    \fi}}
\makeatother

\newcommand{\summary}[1]{%
    \vspace{3pt}
    \noindent
    \large{\textbf{#1}}
}

% Notes lettered a, b, c … on each page (footmisc's perpage). The reset takes effect only from the second TeX pass, so
% on the first the counter runs on through the book, and \alph stops the run past z ("Counter too large": 29 notes in
% the Bible once Mark's were added, 2026-11-16). Past 26 the first pass shows a number instead; the final pass never
% does.
\makeatletter
\renewcommand{\thefootnote}{\ifnum\c@footnote>26 \number\c@footnote\else\@alph\c@footnote\fi}
\makeatother

% set verse margin flush with text
% \setlength{\vleftmargin}{0em}

% default page style
\makepagestyle{plain}
\makeevenhead{plain}{\thepage\hspace{1em}|\hspace{1em}\VersusPrimus}{}{}
\makeoddhead{plain}{}{}{\VersusUltimus\hspace{0.75em}|\hspace{0.75em}\thepage}
\pagestyle{plain}

% Title page
\newlength\drop
\makeatletter
\newcommand*\titleM{\begingroup
\setlength\drop{0.35\textheight}
\centering
\vspace*{\drop}
{\fontsize{35}{45}\bfseries \MakeUppercase{Biblia Sacra}}\\
\vspace*{1em}
\MakeUppercase{\small ex}\\
\vspace*{0.5em}
\MakeUppercase{\LARGE Sēbastiānī Castelliōnis}\\
\vspace*{0.5em}
\MakeUppercase{\large interpretātiōne}
\vfill
\endgroup}
\makeatother


\renewcommand{\contentsname}{\begin{center}\MakeUppercase{ōrdō librōrum}\end{center}}
\AtBeginDocument{\addtocontents{toc}{\protect\thispagestyle{empty}}} 

% A testament's divider: a recto page with its verso blank, so the first book after it opens recto.
\renewcommand{\book}[1]{%
\ClearToRecto
\thispagestyle{empty}%
\beforebookskip
\midpartskip
\begin{center}
\printbooktitle{#1}
\end{center}
\addtocontents{toc}{\begin{center}\Large\bfseries #1\end{center}}%
\ClearToRecto}

% Poetry: no space above or below a poem, so prose and verse share one line grid; no hyphens inside a line.
\makeatletter
\AtBeginEnvironment{verse}{%
  \def\@listi{\leftmargin\leftmargini \topsep\z@ \partopsep\z@ \parsep\z@ \itemsep\z@}%
  \let\@listI\@listi
  \hyphenpenalty=10000 \exhyphenpenalty=10000
  \@rightskip=0pt plus 3em\relax}        % (\list sets \rightskip from it) a long line turns over ragged
\makeatother

\linespread{1.1}  % Standardize line spacing even with superscript verse numbers
